%% ============================================================
%% Introduction to Cognitive Psychology — Week 2 Study Notes
%% ============================================================
\documentclass[11pt,a4paper]{article}

\usepackage[a4paper, left=1.8cm, right=1.8cm, top=1.3cm, bottom=1.8cm]{geometry}
\usepackage{booktabs}
\usepackage{tabularx}
\usepackage{array}
\usepackage{enumitem}
\usepackage{titlesec}
\usepackage{fancyhdr}
\usepackage{fontenc}
\usepackage{inputenc}
\usepackage{microtype}
\usepackage{parskip}
\usepackage{hyperref}
\usepackage{amsmath}
\usepackage{amssymb}

\hypersetup{hidelinks, colorlinks=false}

%% ── Table helpers ─────────────────────────────────────────
\newcolumntype{L}[1]{>{\raggedright\arraybackslash}p{#1}}

%% ── Section headings ──────────────────────────────────────
\titleformat{\section}{\large\bfseries}{}{0em}{}[\titlerule]
\titleformat{\subsection}{\normalsize\bfseries}{}{0em}{}
\titleformat{\subsubsection}{\normalsize\bfseries}{}{0em}{}

\titlespacing{\section}{0pt}{12pt}{4pt}
\titlespacing{\subsection}{0pt}{8pt}{3pt}
\titlespacing{\subsubsection}{0pt}{6pt}{2pt}

%% ── Page style ────────────────────────────────────────────
\pagestyle{fancy}
\fancyhf{}
\renewcommand{\headrulewidth}{0pt}
\fancyfoot[C]{\small Cognitive Psychology --- Week 2 \textbar\ Page \thepage}

%% ── Misc helpers ──────────────────────────────────────────
\newcommand{\bb}[1]{\textbf{#1}}
\setlength{\parskip}{4pt}

%% ═══════════════════════════════════════════════════════════
\begin{document}
\setlength{\arrayrulewidth}{0.4pt}

\begin{center}
\bfseries\LARGE Introduction to Cognitive Psychology\par
\vspace{0.2cm}
\large Assignment 2 Detailed Solutions
\end{center}
\vspace{0.5cm}

%% ════════════════════════════════════════════════════════════
\section*{ASSIGNMENT 2 --- Full Solutions with Explanations}

\subsubsection*{Question 1}
\textit{A toy on the table represents:}
\medskip

\begin{tabular}{L{0.3cm} L{14.9cm}}
& \textbf{A.\ Distal stimulus \checkmark} \\
& B.\ Proximal stimulus \\
& C.\ Retinal image \\
& D.\ Percept \\
\end{tabular}

\vspace{0.3cm}
\noindent\textbf{Ans:} A — Distal Stimulus\\[0.2cm]
\textbf{Why this answer:} \\
A \textbf{distal stimulus} is the actual physical object present in the environment that gives rise to sensory information. The toy sitting on the table is the real-world object — the source of visual input — making it the distal stimulus. This is distinguished from the \textbf{proximal stimulus} (the pattern of energy on the receptor, e.g., the retinal image) and the \textbf{percept} (the brain's interpretation of that input).
\vspace{0.4cm}

\subsubsection*{Question 2}
\textit{A patient with prosopagnosia has difficulty recognizing:}
\medskip

\begin{tabular}{L{0.3cm} L{14.9cm}}
& A.\ Drawing of objects \\
& B.\ Moving objects \\
& \textbf{C.\ Faces \checkmark} \\
& D.\ Sounds \\
\end{tabular}

\vspace{0.3cm}
\noindent\textbf{Ans:} C — Faces\\[0.2cm]
\textbf{Why this answer:} \\
\textbf{Prosopagnosia} is a specific form of \textbf{visual agnosia} in which a patient loses the ability to recognize faces while other forms of object recognition remain relatively intact. This condition demonstrates that face recognition relies on specialized neural mechanisms distinct from general object recognition.
\vspace{0.4cm}

\subsubsection*{Question 3}
\textit{When you let go of a balloon and it floats away from you into the sky, the size of the retinal image gets smaller and smaller. However, you do not perceive the balloon as shrinking. This is an example of:}
\medskip

\begin{tabular}{L{0.3cm} L{14.9cm}}
& A.\ Pattern recognition \\
& B.\ Bottom-up processing \\
& \textbf{C.\ Size constancy \checkmark} \\
& D.\ Figure-ground organization \\
\end{tabular}

\vspace{0.3cm}
\noindent\textbf{Ans:} C — Size Constancy\\[0.2cm]
\textbf{Why this answer:} \\
\textbf{Size constancy} is the perceptual phenomenon whereby the perceived size of an object remains stable even though the size of its \textbf{retinal image} changes as the object moves closer or farther away. In this case, even though the balloon's retinal image shrinks as it floats upward, the brain compensates and you continue to perceive the balloon as being the same actual size.
\vspace{0.4cm}

\subsubsection*{Question 4}
\textit{The meaningful interpretation of a proximal stimulus is called the:}
\medskip

\begin{tabular}{L{0.3cm} L{14.9cm}}
& \textbf{A.\ Percept \checkmark} \\
& B.\ Sensation \\
& C.\ Distal stimulus \\
& D.\ Retinal Image \\
\end{tabular}

\vspace{0.3cm}
\noindent\textbf{Ans:} A — Percept\\[0.2cm]
\textbf{Why this answer:} \\
A \textbf{percept} is the meaningful interpretation that the brain constructs from a \textbf{proximal stimulus} (the raw sensory input reaching the receptors). It represents the final conscious experience of what we perceive — the brain's organized and meaningful understanding of sensory information, going beyond mere sensation.
\vspace{0.4cm}

\subsubsection*{Question 5}
\textit{The two best studied forms of perception are:}
\medskip

\begin{tabular}{L{0.3cm} L{14.9cm}}
& A.\ Visual and Haptic \\
& B.\ Visual and Olfactory \\
& \textbf{C.\ Visual and Auditory \checkmark} \\
& D.\ Auditory and Olfactory \\
\end{tabular}

\vspace{0.3cm}
\noindent\textbf{Ans:} C — Visual and Auditory\\[0.2cm]
\textbf{Why this answer:} \\
\textbf{Vision} and \textbf{audition} (hearing) are the two most extensively studied perceptual systems in cognitive psychology. The vast majority of perception research has focused on how humans process visual and auditory stimuli, due to their dominant role in everyday cognition and behavior.
\vspace{0.4cm}

\subsubsection*{Question 6}
\textit{Which of the following represents a good example of a proximal stimulus?}
\medskip

\begin{tabular}{L{0.3cm} L{14.9cm}}
& A.\ A book on a shelf \\
& B.\ A tree in your yard \\
& C.\ A building on the horizon \\
& \textbf{D.\ The retinal image formed by a tree \checkmark} \\
\end{tabular}

\vspace{0.3cm}
\noindent\textbf{Ans:} D — The retinal image formed by a tree\\[0.2cm]
\textbf{Why this answer:} \\
A \textbf{proximal stimulus} refers to the pattern of stimulation that actually reaches the \textbf{sensory receptors}. The retinal image formed by a tree is the energy pattern (light) projected onto the retina — this is the proximal stimulus. The tree itself (the physical object in the world) is the \textbf{distal stimulus}, not the proximal one.
\vspace{0.4cm}
\newpage

\subsubsection*{Question 7}
\textit{According to Gibson's theory, the acts or behaviors permitted by objects, places, and events are called:}
\medskip

\begin{tabular}{L{0.3cm} L{14.9cm}}
& A.\ Consequences \\
& B.\ Functions \\
& \textbf{C.\ Affordances \checkmark} \\
& D.\ Direct processes \\
\end{tabular}

\vspace{0.3cm}
\noindent\textbf{Ans:} C — Affordances\\[0.2cm]
\textbf{Why this answer:} \\
In James J. \textbf{Gibson's} ecological theory of \textbf{direct perception}, \textbf{affordances} are the actions, behaviors, or interactions that objects, surfaces, places, and events \textit{allow} or \textit{invite}. For example, a chair affords sitting, and a flat surface affords walking. Affordances are perceived directly from the environment without the need for higher-level cognitive processing.
\vspace{0.4cm}

\subsubsection*{Question 8}
\textit{A stencil provides a good analogy for the theory of:}
\medskip

\begin{tabular}{L{0.3cm} L{14.9cm}}
& A.\ Prototype matching \\
& \textbf{B.\ Template matching \checkmark} \\
& C.\ Good continuation \\
& D.\ Featural analysis \\
\end{tabular}

\vspace{0.3cm}
\noindent\textbf{Ans:} B — Template Matching\\[0.2cm]
\textbf{Why this answer:} \\
\textbf{Template matching} theory proposes that pattern recognition occurs by comparing an incoming stimulus to a set of stored, exact internal patterns (templates). A \textbf{stencil} is an excellent analogy because it represents a rigid, predefined shape — recognition happens only when the stimulus matches the template precisely. This contrasts with \textbf{prototype matching}, which uses flexible, averaged representations.
\vspace{0.4cm}

\subsubsection*{Question 9}
\textit{\_\_\_ are to visual perception what phonemes are to language, according to Biederman.}
\medskip

\begin{tabular}{L{0.3cm} L{14.9cm}}
& A.\ Receptors \\
& B.\ Retinas \\
& C.\ Distal stimuli \\
& \textbf{D.\ Geons \checkmark} \\
\end{tabular}

\vspace{0.3cm}
\noindent\textbf{Ans:} D — Geons\\[0.2cm]
\textbf{Why this answer:} \\
According to Irving \textbf{Biederman's Recognition-by-Components (RBC) theory}, \textbf{geons} (geometric ions) are the basic three-dimensional shapes (e.g., cylinders, cones, blocks) that serve as the fundamental building blocks of object recognition. Just as \textbf{phonemes} are the basic units of spoken language, geons are the basic units of visual object perception — a small set of elements that combine to represent a vast number of objects.
\vspace{0.4cm}

\subsubsection*{Question 10}
\textit{The principle governing reversible figures is:}
\medskip

\begin{tabular}{L{0.3cm} L{14.9cm}}
& \textbf{A.\ Figure-ground organization \checkmark} \\
& B.\ Size constancy \\
& C.\ Dimensionality \\
& D.\ Retinal imagery \\
\end{tabular}

\vspace{0.3cm}
\noindent\textbf{Ans:} A — Figure-ground organization\\[0.2cm]
\textbf{Why this answer:} \\
\textbf{Reversible figures} (such as the Rubin vase/faces illusion) are governed by the principle of \textbf{figure-ground organization}. This Gestalt principle describes how perception alternates between two competing interpretations — what is seen as the \textit{figure} (the focal object) versus the \textit{ground} (the background). In reversible figures, the figure and ground swap, creating two distinct percepts from the same image.
\vspace{0.4cm}

\medskip
\subsection*{Assignment Quick Answer Summary}

\begin{center}
\renewcommand{\arraystretch}{1.4}
\begin{tabular}{|L{0.8cm}|L{6.2cm}|L{8.2cm}|}
\hline
\bb{Q\#} & \bb{Topic} & \bb{Correct Answer} \\
\hline
1 & Distal vs. Proximal Stimulus & A — Distal Stimulus \\
2 & Prosopagnosia & C — Faces \\
3 & Perceptual Constancy & C — Size Constancy \\
4 & Percept Definition & A — Percept \\
5 & Most Studied Perception Forms & C — Visual and Auditory \\
6 & Proximal Stimulus Example & D — Retinal Image of a Tree \\
7 & Gibson's Affordances & C — Affordances \\
8 & Template Matching & B — Template Matching \\
9 & Biederman's Geons & D — Geons \\
10 & Reversible Figures & A — Figure-Ground Organization \\
\hline
\end{tabular}
\end{center}

\medskip
\end{document}
